\section*{Appendix B. Lean Formalization}
\addcontentsline{toc}{section}{Appendix B. Lean Formalization}

The file \texttt{lean/RMTChecks.lean} in the replication repository contains 28 theorems that compile under Lean~4.29.0 with the standard distribution alone \citep{demoura2021lean}. It does not depend on mathlib, and it contains no \texttt{sorry} and no \texttt{native\_decide}. The file is checked with \texttt{lean RMTChecks.lean}. The purpose of the formalization is limited and precise. Lean certifies the combinatorial and algebraic steps that connect the analytic arguments of Appendix~A. It does not certify the measure-theoretic statements themselves, namely Wishart moments, the Bartlett decomposition, digamma asymptotics, weak convergence to the Wachter law and Girsanov's theorem. Table~\ref{tab:lean} lists the theorems by group.

\begin{table}[htbp]
\centering
\caption{Theorems in \texttt{lean/RMTChecks.lean}. ``Kernel computation'' means proved by \texttt{decide +kernel}, which evaluates both sides in the Lean kernel with exact integer arithmetic. ``Symbolic'' means proved for all values of the variables, by induction or by the \texttt{grind} decision procedure for commutative rings and fields.}
\label{tab:lean}
\small
\begin{tabular}{>{\raggedright}p{5.4cm}p{7.2cm}l}
\toprule
Theorem & Statement & Method\\
\midrule
\multicolumn{3}{l}{\emph{Moment combinatorics}}\\
\texttt{mp\_moments\_are\_narayana} & recursion \eqref{eq:mp-recursion} gives Narayana polynomials, $k\le10$ & kernel\\
\texttt{mp\_low\_moments} & $m_1,\ldots,m_4$ as quoted in the text & kernel\\
\texttt{catalan\_recursion\_matches} & Catalan recursion equals $\binom{2n}n/(n+1)$, $n\le12$ & kernel\\
\texttt{narayana\_rows\_sum\_}\newline\texttt{to\_catalan} & Narayana rows sum to Catalan numbers, $k\le10$ & kernel\\
\texttt{semicircle\_moments\_catalan} & Wigner moment recursion gives Catalan numbers & kernel\\
\midrule
\multicolumn{3}{l}{\emph{Spurious information}}\\
\texttt{sumTo\_add}, \newline\texttt{sumTo\_sub} & splitting and linearity of finite sums & symbolic\\
\texttt{sumTo\_nonneg}, \newline\texttt{sumTo\_pos} & sign of sums of nonnegative or positive terms & symbolic\\
\texttt{twiceBias\_eq} & telescoping form of \eqref{eq:spurious} & symbolic\\
\texttt{twiceBias\_nonneg}, \texttt{twiceBias\_pos} & positivity for antitone and strictly antitone sequences & symbolic\\
\texttt{twiceBias\_noise}, \texttt{twiceBias\_noise\_pos} & exact increment and strict growth in uninformative coordinates & symbolic\\
\texttt{twiceBias\_symm} & symmetry in $(p,r)$ & symbolic\\
\midrule
\multicolumn{3}{l}{\emph{Phase transitions}}\\
\texttt{mp\_edges} & sum and product of the edges $(1\pm\sqrt q)^2$ & symbolic\\
\texttt{bbp\_outlier\_at\_threshold} & $\lambda(1+\sqrt q)=(1+\sqrt q)^2$ & symbolic\\
\texttt{bbp\_overlap\_at\_threshold} & overlap numerator vanishes at threshold & symbolic\\
\texttt{bbp\_inverse} & outlier satisfies the inversion quadratic & symbolic\\
\texttt{cca\_outlier\_at\_threshold} & $\gamma(t_c)=\omega_+$ in \eqref{eq:cca-spike} & symbolic\\
\texttt{cca\_outlier\_at\_one} & $\gamma(1)=1$ & symbolic\\
\midrule
\multicolumn{3}{l}{\emph{Dynamics, costs and value}}\\
\texttt{mou\_trace\_moment} & $\E\tr H^2=(N+1)/(2a)$ from entry variances & symbolic\\
\texttt{mou\_rate\_closed}, \texttt{mou\_eigen\_closed} & entrywise and eigenvalue rates equal $N(N+1)(b-a)^2/(8a)$ & symbolic\\
\texttt{mou\_eigen\_sufficient} & equality of the two rates & symbolic\\
\texttt{stein\_trace} & $Nn/(n-N-1)-N=N(N+1)/(n-N-1)$ & symbolic\\
\texttt{gmv\_ratio} & product form of the portfolio ratios & symbolic\\
\texttt{cara\_scalar} & scalar determinant and inverse identities & symbolic\\
\bottomrule
\end{tabular}
\end{table}

The central statements about spurious information are proved for an arbitrary sequence $c:\mathbb N\to\mathbb Q$, which plays the role of $i\mapsto\psi((n-i)/2)$. The only property assumed is that the sequence is antitone, so the formal results apply to the digamma terms once the monotonicity of $\psi$ is granted. The two key statements read as follows in the source.

\begin{lstlisting}
def twiceBias (c : Nat → Rat) (p r : Nat) : Rat :=
  sumTo c p + sumTo c r - sumTo c (p + r)

theorem twiceBias_eq (c : Nat → Rat) (p r : Nat) :
    twiceBias c p r = sumTo (fun i => c i - c (p + i)) r

theorem twiceBias_noise_pos (c : Nat → Rat)
    (anti : ∀ i j, i < j → c j < c i) (p r m : Nat)
    (hp : 0 < p) (hm : 0 < m) :
    twiceBias c p r < twiceBias c p (r + m)
\end{lstlisting}

The information-rate result is stated with the stationary entry variances and diffusion coefficients written out explicitly, so that the closed form is checked against the same bookkeeping used in the Girsanov computation:

\begin{lstlisting}
def mouRateEntries (N a b : Rat) : Rat :=
  (1 / 2) * (b - a) ^ 2 *
    (N * (N - 1) / 2 * ((1 / (2 * a * N)) / (1 / N))
      + N * ((1 / (a * N)) / (2 / N)))

theorem mou_eigen_sufficient (N a b : Rat) (hN : N ≠ 0) (ha : a ≠ 0) :
    mouRateEigen N a b = mouRateEntries N a b
\end{lstlisting}
